%% ============================================================
%% MininetGym — AAMAS 2026 Poster
%% A0 Portrait · Compile: pdflatex poster_aamas2026.tex
%% Required packages: beamerposter, tcolorbox, tikz, booktabs
%% Install: sudo apt install texlive-full  (or tlmgr install ...)
%% ============================================================
\documentclass[final]{beamer}
\usepackage[orientation=portrait,size=a0,scale=1.32]{beamerposter}
\usepackage{graphicx}
\usepackage{booktabs}
\usepackage{xcolor}
\usepackage{amsmath}
\usepackage{enumitem}
\usepackage{tikz}
\usetikzlibrary{shapes.geometric,arrows.meta,positioning,fit,backgrounds,calc}
\usepackage{tcolorbox}
\tcbuselibrary{skins,breakable,raster}

%% ---- Color palette ----------------------------------------
\definecolor{cprimary}{HTML}{0D47A1}   % deep blue
\definecolor{csecond}{HTML}{00695C}    % teal
\definecolor{caccent}{HTML}{E64A19}    % orange-red (attacks)
\definecolor{clightbg}{HTML}{F0F4F8}
\definecolor{cdark}{HTML}{212121}
\definecolor{cmid}{HTML}{607D8B}
\definecolor{cred}{HTML}{B71C1C}
\definecolor{cyellow}{HTML}{F57F17}

%% ---- Beamer theme overrides --------------------------------
\setbeamercolor{block title}{fg=white, bg=cprimary}
\setbeamercolor{block body}{fg=cdark, bg=clightbg}
\setbeamercolor{block title alerted}{fg=white, bg=caccent}
\setbeamercolor{block body alerted}{fg=cdark, bg=clightbg}
\setbeamercolor{headline}{fg=white, bg=cprimary!95!black}
\setbeamercolor{footline}{fg=white, bg=cprimary!80!black}
\setbeamercolor{separation line}{bg=caccent}
\setbeamerfont{block title}{size=\large, series=\bfseries}
\setbeamertemplate{navigation symbols}{}
\setbeamertemplate{itemize item}{\color{cprimary}$\blacktriangleright$}

%% ---- tcolorbox styles -------------------------------------
\tcbset{
  scenbox/.style={
    enhanced, arc=5pt, boxrule=1.5pt,
    colback=clightbg, colframe=cprimary,
    colbacktitle=cprimary, coltitle=white,
    fonttitle=\bfseries\normalsize,
    attach boxed title to top left={yshift=-2.5mm, xshift=5mm},
    boxed title style={arc=3pt, boxrule=0.8pt},
    left=4pt, right=4pt, top=6pt, bottom=4pt
  },
  demobox/.style={
    enhanced, arc=5pt, boxrule=1.8pt,
    colback=clightbg, colframe=csecond,
    colbacktitle=csecond, coltitle=white,
    fonttitle=\bfseries\normalsize,
    attach boxed title to top left={yshift=-2.5mm, xshift=5mm},
    boxed title style={arc=3pt, boxrule=0.8pt},
    left=4pt, right=4pt, top=6pt, bottom=4pt
  },
  screenshot/.style={
    enhanced, arc=4pt, boxrule=1pt,
    colback=cmid!12, colframe=cmid!50,
    halign=center, valign=center,
    fontupper=\color{cmid}\itshape\large,
    left=2pt, right=2pt
  }
}

%% ---- Headline ---------------------------------------------
\setbeamertemplate{headline}{%
  \leavevmode
  \begin{beamercolorbox}[wd=\paperwidth]{headline}
    \vspace{1.8cm}
    \begin{columns}[c,totalwidth=\paperwidth]
      \begin{column}{0.10\paperwidth}\centering
        %% Logo placeholder — replace with \includegraphics
        \begin{tikzpicture}
          \node[circle,fill=white!20,minimum size=5.5cm,
                draw=white,line width=3pt,
                font=\bfseries\Large\color{white}] {LOGO};
        \end{tikzpicture}
      \end{column}
      \begin{column}{0.80\paperwidth}\centering
        {\fontsize{68}{74}\selectfont\bfseries\color{white}MininetGym}\\[0.6cm]
        {\fontsize{38}{44}\selectfont\color{white!85!cprimary}%
          A Live Demonstration of RL-Based Cybersecurity Training}\\[0.9cm]
        {\fontsize{28}{34}\selectfont\color{white}%
          Salvo Finistrella \quad$\cdot$\quad [Co-authors]}\\[0.4cm]
        {\fontsize{22}{28}\selectfont\color{white!75!cprimary}%
          DIPI — University of Modena and Reggio Emilia \quad$\cdot$\quad
          AAMAS 2026, Paphos, Cyprus}
      \end{column}
      \begin{column}{0.10\paperwidth}\centering
        %% AAMAS logo placeholder
        \begin{tikzpicture}
          \node[rectangle,fill=white!20,minimum width=5cm,minimum height=5cm,
                draw=white,line width=3pt,
                font=\bfseries\normalsize\color{white},align=center] {AAMAS\\2026};
        \end{tikzpicture}
      \end{column}
    \end{columns}
    \vspace{1.8cm}
  \end{beamercolorbox}
  \begin{beamercolorbox}[wd=\paperwidth,ht=0.4cm]{separation line}\end{beamercolorbox}
}

%% ---- Footline ---------------------------------------------
\setbeamertemplate{footline}{%
  \begin{beamercolorbox}[wd=\paperwidth,ht=3cm,dp=0.5cm]{footline}
    \centering\vspace{0.6cm}
    {\large\color{white}%
      \textbf{AAMAS 2026} $\cdot$ Paphos, Cyprus $\cdot$
      \texttt{github.com/dipi-unimore/mininet-gym} $\cdot$
      \texttt{finix77s@gmail.com} $\cdot$
      Paper: \emph{doi.org/10.1016/j.simpa.2025...}%
    }
  \end{beamercolorbox}
}

%% ===========================================================
\begin{document}
\begin{frame}[t]
\vspace{0.5cm}

\begin{columns}[T,totalwidth=\textwidth]

%% ============================================================
%% COLUMN 1  (29%) — Introduction + Architecture
%% ============================================================
\begin{column}{0.29\textwidth}

% ---- Abstract ----
\begin{block}{Overview}
\large
\textbf{MininetGym} is an open-source, modular platform that turns a live
Software-Defined Network (SDN) emulator into a \textbf{Gymnasium}
environment for training Reinforcement Learning agents on
\textbf{real cybersecurity tasks}.

\medskip
Agents observe live OpenFlow traffic counters, classify hosts as
\emph{normal / victim / attacker}, and can \textbf{trigger SDN
drop-rules} via OpenDayLight to mitigate ongoing DDoS attacks —
all within a reproducible, browser-controlled experiment.
\end{block}

\vspace{0.7cm}

% ---- Architecture TikZ ----
\begin{block}{Framework Architecture}
\centering
\begin{tikzpicture}[
    every node/.style={font=\normalsize},
    box/.style={rectangle, rounded corners=6pt,
                minimum width=13.5cm, minimum height=1.7cm,
                text width=13cm, align=center,
                draw, thick, inner sep=6pt},
    sbox/.style={rectangle, rounded corners=5pt,
                 minimum width=6cm, minimum height=1.5cm,
                 text width=5.7cm, align=center,
                 draw, thick, inner sep=5pt},
    arr/.style={-{Stealth[length=7pt,width=5pt]}, line width=2pt}
  ]

  %% Layer 1 — Web UI
  \node[box, fill=cprimary!18, draw=cprimary] (web) at (0,0)
    {\textbf{Web Dashboard}\\[-2pt]
     \small Flask · SocketIO · Chart.js · Bootstrap};

  %% Layer 2 — Agent Manager
  \node[box, fill=csecond!15, draw=csecond, below=1.1cm of web] (mgr)
    {\textbf{Agent Manager}\\[-2pt]
     \small Training · Evaluation · Metrics · PDF Export};

  %% Layer 3 — RL Agents (two side-by-side)
  \node[sbox, fill=cyellow!15, draw=cyellow!70,
        below=1.2cm of mgr, xshift=-3.8cm] (tabular)
    {\textbf{Tabular Agents}\\[-2pt]
     \small Q-Learning · SARSA};
  \node[sbox, fill=caccent!12, draw=caccent!60,
        below=1.2cm of mgr, xshift=3.8cm] (deep)
    {\textbf{Deep RL (SB3)}\\[-2pt]
     \small DQN · PPO · A2C};

  %% Layer 4 — Gymnasium Environments
  \node[box, fill=cprimary!10, draw=cprimary!60,
        below=3.0cm of mgr] (gym)
    {\textbf{Gymnasium Environments}\\[-2pt]
     \small Classif. · Attack-Net · Attack-PerHost · MARL};

  %% Layer 5 — Mininet
  \node[box, fill=csecond!20, draw=csecond!70,
        below=1.1cm of gym] (mnt)
    {\textbf{Mininet SDN Network}\\[-2pt]
     \small OVS Switch · 5 Hosts · 5 IoT Devices · OpenFlow Monitor};

  %% Layer 6 — SDN Controller + Traffic Gen
  \node[sbox, fill=cred!10, draw=cred!50,
        below=1.2cm of mnt, xshift=-3.8cm] (odl)
    {\textbf{OpenDayLight}\\[-2pt]
     \small SDN Controller · Drop Rules};
  \node[sbox, fill=cred!10, draw=cred!50,
        below=1.2cm of mnt, xshift=3.8cm] (tgen)
    {\textbf{Attack Generator}\\[-2pt]
     \small UDP/TCP/ICMP Flood · Slowloris};

  %% Arrows
  \draw[arr, color=cprimary]  (web.south)  -- (mgr.north);
  \draw[arr, color=csecond]   (mgr.south)  -- ++(0,-.5) -| (tabular.north);
  \draw[arr, color=csecond]   (mgr.south)  -- ++(0,-.5) -| (deep.north);
  \draw[arr, color=cprimary!70] (tabular.south) -- ++(0,-.3) |- (gym.west);
  \draw[arr, color=cprimary!70] (deep.south)    -- ++(0,-.3) |- (gym.east);
  \draw[arr, color=csecond]   (gym.south)  -- (mnt.north);
  \draw[arr, color=cred!70]   (mnt.south)  -- ++(0,-.5) -| (odl.north);
  \draw[arr, color=cred!70]   (mnt.south)  -- ++(0,-.5) -| (tgen.north);
  %% SDN feedback (dashed)
  \draw[{Stealth[length=5pt]}-{Stealth[length=5pt]},
        dashed, line width=1.5pt, color=csecond!60]
    (odl.east) -- ++(1,0) node[above,font=\tiny,color=csecond]{SDN\,loop} -- ++(1.2,0);

\end{tikzpicture}
\end{block}

\vspace{0.7cm}

% ---- Tech Stack ----
\begin{block}{Technology Stack}
\large
\setlength{\tabcolsep}{8pt}
\begin{tabular}{@{}ll@{}}
\toprule
\textbf{Layer}        & \textbf{Technology} \\
\midrule
Network emulation     & Mininet 2.3 + Open vSwitch \\
SDN controller        & OpenDayLight (OpenFlow 1.0) \\
RL framework          & OpenAI Gymnasium + SB3 \\
Deep learning         & PyTorch \\
Web API               & Flask + Flask-SocketIO \\
Frontend              & Bootstrap + Chart.js \\
\bottomrule
\end{tabular}
\end{block}

\end{column}

%% ============================================================
%% COLUMN 2  (1% spacer)
%% ============================================================
\begin{column}{0.01\textwidth}\end{column}

%% ============================================================
%% COLUMN 2  (34%) — Scenarios + Agents + Metrics
%% ============================================================
\begin{column}{0.34\textwidth}

\begin{block}{Cybersecurity Scenarios}
\vspace{0.4cm}

%% Scenario 1
\begin{tcolorbox}[scenbox, title={Scenario 1 \textnormal{—} Traffic Classification}]
\large
Classify live network traffic into \textbf{4 classes}:
\textit{None / Ping / UDP / TCP}.\\[5pt]
\begin{tabular}{@{}ll@{}}
Observation: & $\mathbb{R}^{4}$ — global pkt/byte aggregates \\
Action:      & $\{0,1,2,3\}$ \\
Reward:      & graded by distance from correct label \\
\end{tabular}
\end{tcolorbox}

\vspace{0.5cm}

%% Scenario 2
\begin{tcolorbox}[scenbox, title={Scenario 2 \textnormal{—} Network-Level Attack Detection}]
\large
Binary detection of ongoing flooding attacks at the
\textbf{global network level}.\\[5pt]
\begin{tabular}{@{}ll@{}}
Observation: & $\mathbb{R}^{4}$ — pkts, $\Delta$pkts\%, bytes, $\Delta$bytes\% \\
Action:      & \{Normal, Attack\} \\
Defense:     & alert only \\
\end{tabular}
\end{tcolorbox}

\vspace{0.5cm}

%% Scenario 3
\begin{tcolorbox}[scenbox,
  title={Scenario 3 \textnormal{—} Per-Host Detection \textcolor{cyellow}{\normalfont(primary scenario)}}]
\large
Each host is classified as \textbf{Normal / Victim / Attacker}.
Detection of an \emph{attacker} triggers an
\textbf{\color{cred}OpenFlow drop rule} via OpenDayLight.\\[5pt]
\begin{tabular}{@{}ll@{}}
Observation: & $\mathbb{R}^{8}$ per host (TX+RX pkt/byte+deltas) \\
Action:      & \{Normal, Incoming attack, Outgoing attack\} \\
Defense:     & \textcolor{cred}{\textbf{SDN link blocking}} \\
Wrapper:     & \textit{PerHostScanWrapper} (const obs size) \\
\end{tabular}\\[5pt]
\small
Reward asymmetry compensates class imbalance:
$r_{\text{scale}} = \min\!\left(\frac{N-1}{2},\,8\right)$, penalises
false negatives $\times 2$.
\end{tcolorbox}

\vspace{0.5cm}

%% Scenario 4
\begin{tcolorbox}[scenbox,
  title={Scenario 4 \textnormal{—} Multi-Agent Hierarchical (MARL)}]
\large
A \textbf{Coordinator} monitors global traffic and broadcasts alerts;
independent \textbf{Host agents} act on local per-host observations.\\[5pt]
\begin{tabular}{@{}ll@{}}
Coordinator obs: & $\mathbb{R}^{5}$ + message count \\
Host obs:        & $\mathbb{R}^{9}$ + message count \\
Communication:   & shared \textit{InstantState} message bus \\
Training:        & parallel threads (one per agent) \\
\end{tabular}
\end{tcolorbox}

\vspace{0.3cm}
\end{block}

\vspace{0.6cm}

\begin{block}{Supported RL Agents}
\large
\begin{columns}[T]
\begin{column}{0.46\textwidth}
\textbf{Tabular}
\begin{itemize}[leftmargin=1.2em,itemsep=2pt]
  \item Q-Learning\\\small(log-bin discretization)
  \item SARSA (on-policy)
\end{itemize}
\bigskip
\textbf{Supervised Baseline}
\begin{itemize}[leftmargin=1.2em,itemsep=2pt]
  \item Supervised Agent
\end{itemize}
\end{column}
\begin{column}{0.46\textwidth}
\textbf{Deep RL — Stable Baselines 3}
\begin{itemize}[leftmargin=1.2em,itemsep=2pt]
  \item DQN (Deep Q-Network)
  \item PPO (Proximal Policy Opt.)
  \item A2C (Advantage Actor-Critic)
\end{itemize}
\end{column}
\end{columns}
\end{block}

\vspace{0.6cm}

\begin{block}{Key Evaluation Metrics}
\large
\begin{itemize}[leftmargin=1.2em,itemsep=4pt]
  \item \textbf{Accuracy · Precision · Recall · F1} per agent
  \item \textbf{Mitigation Ratio} — fraction of attacking hosts
        successfully blocked by SDN
  \item \textbf{False Negative Rate} — missed attacks
        (high-cost errors, penalised $\times$2)
  \item Live \textbf{reward curves} \& \textbf{confusion matrices}
  \item \textbf{Attack latency} — steps from attack start to blocking
\end{itemize}
\end{block}

\end{column}

%% ============================================================
%% COLUMN 3  (1% spacer)
%% ============================================================
\begin{column}{0.01\textwidth}\end{column}

%% ============================================================
%% COLUMN 3  (35%) — Live Demo Walkthrough
%% ============================================================
\begin{column}{0.35\textwidth}

\begin{block}{Live Demo Walkthrough}
\vspace{0.5cm}

%% ---- STEP 1 ----
\begin{tcolorbox}[demobox,
  title={\large\raisebox{1pt}{\textbf{\textsf{①}}}~Introduction}]
\large
The attendee opens a \textbf{URL} on their device (laptop or mobile).
The single-page application presents three tabs:
\textit{Configuration}, \textit{Training}, and \textit{Results}.

\vspace{0.5cm}
\begin{tcolorbox}[screenshot, height=9cm]
\textbf{[ Screenshot — Landing Page ]}\\[4pt]
Tab navigation · Status bar · System info
\end{tcolorbox}
\end{tcolorbox}

\vspace{0.55cm}

%% ---- STEP 2 ----
\begin{tcolorbox}[demobox,
  title={\large\raisebox{1pt}{\textbf{\textsf{②}}}~Configuration}]
\large
The demonstrator sets the \textbf{network topology}
(5 hosts + 5 IoT), selects \emph{Scenario 3 — attacks\_ho},
adds three agents (PPO, DQN, Q-Learning), and tunes
hyperparameters — all via the web form.

\vspace{0.5cm}
\begin{tcolorbox}[screenshot, height=9cm]
\textbf{[ Screenshot — Configuration Panel ]}\\[4pt]
Topology · Agent cards · Attack parameters
\end{tcolorbox}
\end{tcolorbox}

\vspace{0.55cm}

%% ---- STEP 3 ----
\begin{tcolorbox}[demobox,
  title={\large\raisebox{1pt}{\textbf{\textsf{③}}}~Experiment \textnormal{(also on mobile)}}]
\large
Training starts. \textbf{Reward and accuracy charts} update in
real time via WebSocket. Host-status indicators show
which hosts are under attack and whether the agent has
triggered \textbf{\color{cred}SDN blocking}.
The UI is \textbf{fully responsive} — accessible from a phone.

\vspace{0.4cm}
\begin{columns}[T]
  \begin{column}{0.61\textwidth}
    \begin{tcolorbox}[screenshot, height=12cm]
      \textbf{[ Screenshot — Training Dashboard ]}\\[4pt]
      Live charts · Host status grid
    \end{tcolorbox}
  \end{column}
  \hfill
  \begin{column}{0.36\textwidth}
    \begin{tcolorbox}[screenshot, height=12cm]
      \textbf{[ Mobile View ]}\\[4pt]
      Responsive layout
    \end{tcolorbox}
  \end{column}
\end{columns}
\end{tcolorbox}

\vspace{0.55cm}

%% ---- STEP 4 ----
\begin{tcolorbox}[demobox,
  title={\large\raisebox{1pt}{\textbf{\textsf{④}}}~Result Evaluation}]
\large
The \textit{Results} tab displays per-agent
accuracy tables, confusion matrices, and a
\textbf{mitigation gauge} showing the fraction
of attacks successfully blocked by OpenFlow rules.
A \textbf{PDF report} can be exported in one click.

\vspace{0.4cm}
\begin{columns}[T]
  \begin{column}{0.48\textwidth}
    \begin{tcolorbox}[screenshot, height=10cm]
      \textbf{[ Screenshot — Metrics Table ]}\\[4pt]
      Accuracy · Recall · F1
    \end{tcolorbox}
  \end{column}
  \hfill
  \begin{column}{0.48\textwidth}
    \begin{tcolorbox}[screenshot, height=10cm]
      \textbf{[ Screenshot — Mitigation Gauge ]}\\[4pt]
      SDN block rate
    \end{tcolorbox}
  \end{column}
\end{columns}
\end{tcolorbox}

\end{block}

\end{column}

\end{columns}

\end{frame}
\end{document}
